\chapter*{Conventions}
\addcontentsline{toc}{chapter}{Conventions}

This chapter describes the conventions used throughout the thesis to ensure consistency and clarity.

%%%%%%%%%%%%%%%%%%%%%%%%%%%%%%%%%%%%%%%%%%%%%%%%%%%%%%%%%%%%%%%%%%%%%%%%
\section*{Language}

This thesis is written in British/American English.\footnote{Choose one and be consistent throughout.} Technical terms are introduced in \emph{italics} when first used and defined in the Glossary.

%%%%%%%%%%%%%%%%%%%%%%%%%%%%%%%%%%%%%%%%%%%%%%%%%%%%%%%%%%%%%%%%%%%%%%%%
\section*{Notation}

Describe any mathematical or typographical conventions you follow. For example:

\begin{itemize}
  \item Vectors are denoted by bold lowercase letters, e.g.\ $\mathbf{v}$.
  \item Matrices are denoted by bold uppercase letters, e.g.\ $\mathbf{M}$.
  \item Sets are denoted by calligraphic uppercase letters, e.g.\ $\mathcal{S}$.
  \item Scalar variables are denoted by italic letters, e.g.\ $x$.
\end{itemize}

%%%%%%%%%%%%%%%%%%%%%%%%%%%%%%%%%%%%%%%%%%%%%%%%%%%%%%%%%%%%%%%%%%%%%%%%
\section*{References}

Figures, tables, chapters, and sections are referred to with a capital letter, e.g.\ Figure~\ref{fig:design}, Table~\ref{tab:results}, Chapter~\ref{ch:introduction}, Section~\ref{sec:motivation}.
